\documentclass[12pt, a4paper]{article}

\usepackage[utf8]{inputenc}
\usepackage[T2A]{fontenc}
\usepackage[russian]{babel}
\usepackage{amsmath, amsfonts, amssymb, amsthm}
\usepackage{geometry}
\usepackage{listings}
\usepackage{xcolor}
\usepackage{hyperref}

\geometry{left=25mm, right=15mm, top=20mm, bottom=20mm}

\newtheorem{theorem}{Теорема}
\newtheorem{definition}{Определение}
\newtheorem{lemma}{Лемма}

\definecolor{codegreen}{rgb}{0,0.6,0}
\definecolor{codegray}{rgb}{0.5,0.5,0.5}
\definecolor{codepurple}{rgb}{0.58,0,0.82}
\definecolor{backcolour}{rgb}{0.95,0.95,0.92}

\lstdefinestyle{mystyle}{
    backgroundcolor=\color{backcolour},
    commentstyle=\color{codegreen},
    keywordstyle=\color{blue},
    numberstyle=\tiny\color{codegray},
    stringstyle=\color{codepurple},
    basicstyle=\ttfamily\footnotesize,
    breakatwhitespace=false,
    breaklines=true,
    captionpos=b,
    keepspaces=true,
    numbers=left,
    numbersep=5pt,
    showspaces=false,
    showstringspaces=false,
    showtabs=false,
    tabsize=2,
    literate=
      {А}{{\CYRA}}1
      {Б}{{\CYRB}}1
      {В}{{\CYRV}}1
      {Г}{{\CYRG}}1
      {Д}{{\CYRD}}1
      {Е}{{\CYRE}}1
      {Ё}{{\CYRYO}}1
      {Ж}{{\CYRZH}}1
      {З}{{\CYRZ}}1
      {И}{{\CYRI}}1
      {Й}{{\CYRISHRT}}1
      {К}{{\CYRK}}1
      {Л}{{\CYRL}}1
      {М}{{\CYRM}}1
      {Н}{{\CYRN}}1
      {О}{{\CYRO}}1
      {П}{{\CYRP}}1
      {Р}{{\CYRR}}1
      {С}{{\CYRS}}1
      {Т}{{\CYRT}}1
      {У}{{\CYRU}}1
      {Ф}{{\CYRF}}1
      {Х}{{\CYRH}}1
      {Ц}{{\CYRC}}1
      {Ч}{{\CYRCH}}1
      {Ш}{{\CYRSH}}1
      {Щ}{{\CYRSHCH}}1
      {Ъ}{{\CYRHRDSN}}1
      {Ы}{{\CYRERY}}1
      {Ь}{{\CYRSFTSN}}1
      {Э}{{\CYREREV}}1
      {Ю}{{\CYRYU}}1
      {Я}{{\CYRYA}}1
      {а}{{\cyra}}1
      {б}{{\cyrb}}1
      {в}{{\cyrv}}1
      {г}{{\cyrg}}1
      {д}{{\cyrd}}1
      {е}{{\cyre}}1
      {ё}{{\cyryo}}1
      {ж}{{\cyrzh}}1
      {з}{{\cyrz}}1
      {и}{{\cyri}}1
      {й}{{\cyrishrt}}1
      {к}{{\cyrk}}1
      {л}{{\cyrl}}1
      {м}{{\cyrm}}1
      {н}{{\cyrn}}1
      {о}{{\cyro}}1
      {п}{{\cyrp}}1
      {р}{{\cyrr}}1
      {с}{{\cyrs}}1
      {т}{{\cyrt}}1
      {у}{{\cyru}}1
      {ф}{{\cyrf}}1
      {х}{{\cyrh}}1
      {ц}{{\cyrc}}1
      {ч}{{\cyrch}}1
      {ш}{{\cyrsh}}1
      {щ}{{\cyrshch}}1
      {ъ}{{\cyrhrdsn}}1
      {ы}{{\cyrery}}1
      {ь}{{\cyrsftsn}}1
      {э}{{\cyrerev}}1
      {ю}{{\cyryu}}1
      {я}{{\cyrya}}1
}
\lstset{style=mystyle}

\title{Билет №79 \\ \small Стандарты языков SQL. Интерактивный, встроенный, динамический SQL. (ИТМО) \\ Стандарты языков SQL. Интерактивный, встроенный, динамический SQL. (СпБГУ)}
\author{Сериков Владислав Александрович}
\date{16 мая 2026}

\begin{document}

\maketitle
\tableofcontents
\newpage

\section{Введение: Реляционная модель и SQL}

Язык SQL (Structured Query Language) является фактическим стандартом для работы с реляционными базами данных. Для понимания выразительной мощности SQL необходимо обратиться к теоретическому фундаменту — реляционной модели данных Э. Кодда.

\begin{definition}
\textbf{Реляционная алгебра} — это процедурный язык запросов, состоящий из множества операций, которые принимают в качестве входа одно или два отношения и возвращают в качестве результата новое отношение.
\end{definition}

\begin{definition}
\textbf{Реляционное исчисление} (кортежей или доменов) — это непроцедурный язык запросов, основанный на логике предикатов первого порядка, описывающий, \textit{какие} данные нужно получить, а не \textit{как} их получить.
\end{definition}

Фундаментальным результатом теории баз данных, определяющим требования к полноте языка SQL, является теорема Кодда.

\begin{theorem}[Теорема Кодда о реляционной полноте]
Выразительная мощность реляционной алгебры и безопасного реляционного исчисления совпадает. То есть, любой запрос, который может быть выражен с помощью безопасной формулы реляционного исчисления, может быть выражен средствами реляционной алгебры, и наоборот.
\end{theorem}

\begin{proof}
Доказательство состоит из двух частей.

\textbf{Часть 1. От алгебры к исчислению.}
Покажем, что каждая из пяти базовых операций реляционной алгебры (выборка $\sigma$, проекция $\pi$, декартово произведение $\times$, объединение $\cup$, разность $\setminus$) может быть выражена в исчислении кортежей. Пусть $R$ и $S$ — отношения.
\begin{itemize}
    \item \textit{Выборка} $\sigma_F(R) \equiv \{ t \mid R(t) \land F(t) \}$
    \item \textit{Проекция} $\pi_{A}(R) \equiv \{ t[A] \mid \exists u (R(u) \land t[A] = u[A]) \}$
    \item \textit{Декартово произведение} $R \times S \equiv \{ t \mid \exists u \exists v (R(u) \land S(v) \land t = u \circ v) \}$
    \item \textit{Объединение} $R \cup S \equiv \{ t \mid R(t) \lor S(t) \}$
    \item \textit{Разность} $R \setminus S \equiv \{ t \mid R(t) \land \neg S(t) \}$
\end{itemize}
Таким образом, любое выражение реляционной алгебры может быть сведено к формуле реляционного исчисления.

\textbf{Часть 2. От исчисления к алгебре (структурная индукция).}
Докажем, что для любой безопасной формулы исчисления кортежей $\psi$ существует эквивалентное выражение реляционной алгебры $E$. Доказательство проводится индукцией по построению формулы.

\textit{База индукции:} Формула $\psi$ является атомарной $R(t)$. Эквивалентное алгебраическое выражение — само отношение $R$.

\textit{Индукционный переход:} Предположим, что для безопасных подформул $\psi_1$ и $\psi_2$ существуют алгебраические выражения $E_1$ и $E_2$.
\begin{enumerate}
    \item Если $\psi = \psi_1 \land \psi_2$, то результатом является естественное соединение (или декартово произведение с последующей выборкой): $E_1 \bowtie E_2$.
    \item Если $\psi = \psi_1 \lor \psi_2$, то результатом является объединение: $E_1 \cup E_2$ (с учетом совместимости схем).
    \item Если $\psi = \psi_1 \land \neg \psi_2$ (обратите внимание, $\neg \psi_2$ отдельно небезопасна, но безопасна в конъюнкции), результатом является разность: $E_1 \setminus E_2$.
    \item Если $\psi = \exists x (\psi_1(x))$, где переменная $x$ ограничивается доменом атрибутов, то операция эквивалентна проекции выражения $E_1$, исключающей атрибуты, связанные с $x$: $\pi_{Attr(E_1) \setminus \{x\}}(E_1)$.
\end{enumerate}
Универсальный квантор $\forall$ выражается через $\exists$ и отрицание. Таким образом, по принципу математической индукции, любая безопасная формула транслируется в реляционную алгебру. Теорема доказана.
\end{proof}

Язык SQL является \textbf{реляционно-полным}, так как реализует все операции реляционной алгебры.

\section{Стандарты языка SQL}

С момента своего появления (изначально под названием SEQUEL) язык прошел долгий путь стандартизации комитетами ANSI и ISO/IEC (ISO/IEC 9075).

\begin{itemize}
    \item \textbf{SQL-86 (SQL-87):} Первый официальный стандарт. Описаны базовые конструкции DML (Data Manipulation Language) и DDL (Data Definition Language).
    \item \textbf{SQL-89:} Незначительные доработки, добавлена поддержка ограничения целостности (первичные и внешние ключи).
    \item \textbf{SQL-92 (SQL2):} Крупное обновление. Введены типы данных (VARCHAR, DATE, TIME), стандартизированы операции JOIN, добавлены уровни изоляции транзакций. Стандарт разделен на уровни соответствия (Entry, Intermediate, Full).
    \item \textbf{SQL:1999 (SQL3):} Революционный стандарт. Добавлены регулярные выражения, рекурсивные CTE (Common Table Expressions), триггеры, а главное — объектно-реляционные возможности (пользовательские типы данных, массивы).
    \item \textbf{SQL:2003:} Введены оконные функции (Window functions), последовательности (SEQUENCE) и поддержка XML.
    \item \textbf{SQL:2008:} Улучшена интеграция с XML, добавлен оператор TRUNCATE, расширены оконные функции и триггеры.
    \item \textbf{SQL:2011:} Введена поддержка темпоральных баз данных (периоды действия строк: \texttt{VALID TIME}, \texttt{TRANSACTION TIME}).
    \item \textbf{SQL:2016:} Добавлена встроенная поддержка формата JSON, полиморфные табличные функции (PTF), Row-level security.
    \item \textbf{SQL:2023:} Последний актуальный стандарт. Введение языка Property Graph Queries (SQL/PGQ), позволяющего выполнять графовые запросы к реляционным данным.
\end{itemize}

\section{Интерактивный SQL}

\begin{definition}
\textbf{Интерактивный SQL} — это режим работы, при котором пользователь или администратор вводит SQL-запросы напрямую через интерфейс командной строки (CLI) или графический интерфейс (GUI), и СУБД немедленно их выполняет, возвращая результат на экран.
\end{definition}

\textbf{Особенности:}
\begin{itemize}
    \item Запросы независимы от контекста языка программирования.
    \item Используется для администрирования, ad-hoc запросов (запросов по требованию) и аналитики.
    \item Результат обычно представлен в виде таблицы на экране (или экспортируется в файл).
\end{itemize}

\section{Встроенный SQL (Embedded SQL)}

Интерактивный SQL не подходит для создания сложных прикладных программ, так как не обладает полнотой по Тьюрингу (в нем нет полноценных циклов, ветвлений и т.д., за исключением расширений вроде PL/pgSQL). Поэтому SQL встраивают в процедурные языки программирования (C, Java, COBOL) — так называемые \textit{хост-языки}.

\textbf{Проблема разрыва импеданса (Impedance Mismatch):}
Реляционная модель оперирует множествами (отношениями), в то время как традиционные языки программирования оперируют отдельными скалярными переменными и записями. Решением этой проблемы является механизм \textbf{курсоров}.

\begin{definition}
\textbf{Встроенный (статический) SQL} — это технология, при которой SQL-инструкции жестко закодированы (hardcoded) в исходном коде программы на хост-языке и обрабатываются специальным прекомпилятором до основной компиляции.
\end{definition}

Во встроенном SQL используются конструкции \texttt{EXEC SQL}, а для обмена данными между SQL и программой применяются \textit{хост-переменные} (обычно предваряются двоеточием \texttt{:}).

\subsection{Алгоритм работы с курсором во встроенном SQL}

Чтобы обработать набор строк, возвращаемый запросом, используется курсор.

\textbf{Шаги алгоритма:}
\begin{enumerate}
    \item \textbf{Объявление хост-переменных:} Выделение памяти в программе для хранения извлекаемых данных.
    \item \textbf{Объявление курсора (\texttt{DECLARE}):} Связывание имени курсора с конкретным \texttt{SELECT}-запросом. На этом этапе запрос не выполняется.
    \item \textbf{Открытие курсора (\texttt{OPEN}):} СУБД выполняет запрос, формирует результирующий набор данных и устанавливает указатель перед первой строкой.
    \item \textbf{Извлечение данных (\texttt{FETCH}):} В цикле указатель сдвигается на следующую строку. Данные текущей строки копируются в хост-переменные.
    \item \textbf{Проверка статуса (SQLCODE / SQLSTATE):} Проверка на достижение конца результирующего множества (например, \texttt{SQLSTATE = '02000'}).
    \item \textbf{Закрытие курсора (\texttt{CLOSE}):} Освобождение ресурсов, выделенных под результирующее множество.
\end{enumerate}

\textbf{Минимальный пример (C + Embedded SQL):}
\begin{lstlisting}[language=C]
EXEC SQL BEGIN DECLARE SECTION;
    int emp_id;
    char emp_name[50];
EXEC SQL END DECLARE SECTION;

// Шаг 2. Объявление курсора
EXEC SQL DECLARE emp_cursor CURSOR FOR
    SELECT id, name FROM Employees WHERE department = 'IT';

// Шаг 3. Открытие курсора
EXEC SQL OPEN emp_cursor;

// Шаг 4-5. Извлечение и проверка в цикле
while (1) {
    EXEC SQL FETCH emp_cursor INTO :emp_id, :emp_name;

    // Проверка на отсутствие данных (Конец множества)
    if (sqlca.sqlcode == 100) break;

    printf("ID: %d, Name: %s\n", emp_id, emp_name);
}

// Шаг 6. Закрытие курсора
EXEC SQL CLOSE emp_cursor;
\end{lstlisting}

\section{Динамический SQL (Dynamic SQL)}

Во встроенном статическом SQL структура запроса (таблицы, столбцы) должна быть известна на момент компиляции. Если структура запроса формируется пользователем во время работы программы (например, конструктор отчетов), статический SQL неприменим.

\begin{definition}
\textbf{Динамический SQL} — это метод работы, при котором SQL-запросы формируются в виде строковых переменных во время выполнения программы (runtime) и передаются СУБД для компиляции и исполнения.
\end{definition}

Для обработки динамических запросов используются операторы \texttt{PREPARE} (подготовка) и \texttt{EXECUTE} (исполнение).

\subsection{Алгоритм выполнения параметризованного динамического запроса}

\textbf{Шаги алгоритма:}
\begin{enumerate}
    \item \textbf{Формирование строки запроса:} Программа генерирует строку, содержащую SQL-код. Вместо литералов (конкретных значений) рекомендуется использовать параметры-заполнители (обычно \texttt{?}), чтобы избежать SQL-инъекций и оптимизировать кеш планов запросов.
    \item \textbf{Подготовка (\texttt{PREPARE}):} Строка отправляется в СУБД. СУБД проводит синтаксический анализ, семантическую проверку и строит план выполнения. Готовому запросу присваивается идентификатор (имя).
    \item \textbf{Выполнение (\texttt{EXECUTE}):} СУБД выполняет подготовленный план, подставляя на место заполнителей значения из хост-переменных, переданных через клаузулу \texttt{USING}.
    \item \textbf{Освобождение (\texttt{DEALLOCATE}):} Удаление подготовленного запроса из памяти СУБД.
\end{enumerate}

\textbf{Минимальный пример (Псевдокод / C):}
\begin{lstlisting}[language=C]
EXEC SQL BEGIN DECLARE SECTION;
    char query_string[] = "UPDATE Employees SET salary = salary * ? WHERE department = ?";
    float multiplier = 1.1;
    char dept[] = "Sales";
EXEC SQL END DECLARE SECTION;

// Шаг 2. Подготовка
EXEC SQL PREPARE stmt FROM :query_string;

// Шаг 3. Выполнение с передачей параметров
EXEC SQL EXECUTE stmt USING :multiplier, :dept;

// Шаг 4. Освобождение ресурсов
EXEC SQL DEALLOCATE PREPARE stmt;
\end{lstlisting}

\textit{Примечание:} Если динамический запрос является оператором \texttt{SELECT}, возвращающим множество строк, то после шага \texttt{PREPARE} необходимо использовать динамический курсор (объявляется для подготовленного выражения \texttt{stmt}, а не для статической строки), после чего применять алгоритм обхода курсора из Раздела 4.1. Для однократных запросов без параметров можно использовать сокращенную форму: \texttt{EXECUTE IMMEDIATE :query\_string}.

\end{document}
